\documentclass[11pt]{article}

\usepackage[margin=1in]{geometry}
\usepackage{amsmath,amssymb,amsthm,mathtools,bm}
\usepackage{enumitem}
\usepackage{microtype}
\usepackage[hidelinks]{hyperref}

\newtheorem{theorem}{Theorem}[section]
\newtheorem{proposition}[theorem]{Proposition}
\newtheorem{lemma}[theorem]{Lemma}
\newtheorem{corollary}[theorem]{Corollary}
\newtheorem{remark}[theorem]{Remark}

\newcommand{\R}{\mathbb{R}}
\newcommand{\E}{\mathbb{E}}
\newcommand{\Pp}{\mathbb{P}}
\newcommand{\1}{\mathbf{1}}
\newcommand{\op}{\mathrm{op}}
\newcommand{\Tr}{\operatorname{Tr}}
\newcommand{\Diag}{\operatorname{Diag}}
\newcommand{\softmax}{\operatorname{softmax}}
\newcommand{\Var}{\operatorname{Var}}
\newcommand{\ip}[2]{\left\langle #1,#2\right\rangle_F}
\newcommand{\norm}[1]{\left\|#1\right\|}
\newcommand{\abs}[1]{\left|#1\right|}
\newcommand{\set}[1]{\left\{#1\right\}}
\newcommand{\brak}[1]{\left[#1\right]}
\newcommand{\paren}[1]{\left(#1\right)}
\newcommand{\sphere}{\mathbb{S}}

\title{A Finite-Dimensional Technical Note on Global Smoothness of\\
Gaussian Associative-Memory Cross-Entropy\\
under Frobenius and Operator-Norm Geometries}
\author{}
\date{}

\begin{document}
\maketitle

\begin{abstract}
We give a self-contained finite-dimensional proof of global smoothness bounds for a
Gaussian associative-memory softmax loss.  The proof is organized around one explicit
good event, and every failure probability is stated in its original exponential-tail
form; no confidence parameter is introduced.  For Zipf weights
\[
 p_i=\frac{i^{-\alpha}}{H_{N,\alpha}},
 \qquad
 H_{n,\alpha}=\sum_{k=1}^n k^{-\alpha},
 \qquad \alpha\ge 0,
\]
we prove, simultaneously with an explicit probability, that
\[
 \beta_F\asymp \frac1d+p_1,
 \qquad
 \beta_{\op}\gtrsim \frac1d+P_{\min\{d,N\}},
\]
and we establish an explicit finite-dimensional upper bound for
\(\beta_{\op}\).  In the regime \(N=d^{\Theta(1)}\), the operator-norm
bound matches the lower bound up to polylogarithmic factors.
\end{abstract}

\tableofcontents

\section{Setup, definitions, and main theorem}
\label{sec:setup}

\subsection{Random model and loss}

For an integer \(n\ge 1\), write \([n]:=\{1,\ldots,n\}\).  Fix integers
\(d,N\ge 2\).  Let
\[
 u_1,\ldots,u_N,\ v_1,\ldots,v_N
\]
be mutually independent random vectors with
\[
 u_j\sim \mathcal N\!\left(0,\frac1d I_d\right),
 \qquad
 v_i\sim \mathcal N\!\left(0,\frac1d I_d\right).
\]
The vectors \(u_j\) are the output embeddings and the vectors \(v_i\) are the
input embeddings.  Put
\[
 U:=[u_1,\ldots,u_N]\in\R^{d\times N}.
\]

Fix a Zipf exponent \(\alpha\ge 0\), and define
\[
 H_{n,\alpha}:=\sum_{k=1}^n k^{-\alpha},
 \qquad
 p_i:=\frac{i^{-\alpha}}{H_{N,\alpha}},
 \qquad i\in[N].
\]
Thus \(p_1\ge\cdots\ge p_N>0\) and \(\sum_{i=1}^N p_i=1\).  For
\(s\in[N]\), define the head mass
\[
 P_s:=\sum_{i=1}^s p_i=\frac{H_{s,\alpha}}{H_{N,\alpha}}.
\]

For \(W\in\R^{d\times d}\), define the logits and softmax probabilities
\[
 z_{ij}(W):=u_j^\top Wv_i,
 \qquad
 q_{ij}(W):=\frac{e^{z_{ij}(W)}}{\sum_{k=1}^N e^{z_{ik}(W)}}.
\]
The loss associated with input \(i\) and the population loss are
\[
 \ell_i(W)
 :=
 \log\!\left(\sum_{j=1}^N e^{u_j^\top Wv_i}\right)-u_i^\top Wv_i,
 \qquad
 L(W):=\sum_{i=1}^N p_i\ell_i(W).
\]
The target term \(-u_i^\top Wv_i\) is linear in \(W\), and therefore it does
not appear in the Hessian.

We use the Frobenius inner product
\[
 \ip{A}{B}:=\Tr(A^\top B).
\]
The Frobenius and operator norms are denoted by \(\norm{\cdot}_F\) and
\(\norm{\cdot}_{\op}\), respectively.  Define the global quadratic
smoothness constants
\[
 \beta_F
 :=
 \sup_{W\in\R^{d\times d}}
 \sup_{H\ne 0}
 \frac{\nabla^2L(W)[H,H]}{\norm{H}_F^2},
\]
\[
 \beta_{\op}
 :=
 \sup_{W\in\R^{d\times d}}
 \sup_{H\ne 0}
 \frac{\nabla^2L(W)[H,H]}{\norm{H}_{\op}^2}.
\]
Section~\ref{sec:hessian} proves that these are exactly the global Lipschitz
constants of the gradient from the corresponding primal norm to its dual norm.

\subsection{Parameters and explicit failure probability}

Set
\[
 m:=\min\{d,N\},
 \qquad
 K:=\left\lfloor\frac{m}{256}\right\rfloor.
\]
When \(N>d\), define
\[
 J:=\left\lceil\log_2\frac Nd\right\rceil,
 \qquad
 \Xi_{\alpha,J}:=\sum_{r=0}^{J-1}2^{-\alpha r}.
\]
When \(N\le d\), set \(J:=0\) and \(\Xi_{\alpha,J}:=0\).  In all cases, set
\[
 L_N:=\log(2N^2),
 \qquad
 T_d:=d^2\log(3d)+d.
\]

The selector construction below has failure probability
\[
 \varepsilon_{\mathrm{sel}}(d,N,K)
 :=
 4K e^{-d/32}
 +2K e^{-d/48}
 +K(N-2)e^{-d/64}.
\]
Define the total explicit failure bound
\begin{align}
 \varepsilon_{\mathrm{tot}}(d,N)
 :=\;&
 2N e^{-d/2}
 +e^{-d}
 +e^{-d/16}
 +\varepsilon_{\mathrm{sel}}(d,N,K)
 \notag\\
 &\quad
 +4\cdot 17^K e^{-d/32}
 +J e^{-d}.
 \label{eq:total-failure}
\end{align}
No independence between the events contributing to
\eqref{eq:total-failure} is asserted or needed; the final estimate follows from
a union bound and one conditional union bound for the dyadic blocks.

\begin{theorem}[Finite-dimensional global smoothness]
\label{thm:main}
Assume
\[
 \alpha\ge 0,
 \qquad
 m=\min\{d,N\}\ge 512.
\]
Then, with probability at least
\[
 \boxed{1-\varepsilon_{\mathrm{tot}}(d,N)},
\]
all of the following conclusions hold simultaneously.

\paragraph{Frobenius geometry.}
\[
 \boxed{
 \frac1{64}\left(\frac1d+p_1\right)
 \le \beta_F
 \le 80\left(\frac1d+p_1\right).
 }
\]
Equivalently, because \(p_1=H_{N,\alpha}^{-1}\),
\[
 \frac1{64}\left(\frac1d+\frac1{H_{N,\alpha}}\right)
 \le \beta_F
 \le 80\left(\frac1d+\frac1{H_{N,\alpha}}\right).
\]

\paragraph{Operator geometry: lower bound.}
\[
 \boxed{
 \beta_{\op}
 \ge
 \frac1{294912}
 \left(\frac1d+P_m\right).
 }
\]

\paragraph{Operator geometry: upper bound.}
In every case,
\[
 \boxed{\beta_{\op}\le 16.}
\]
If \(N>d\), then the sharper estimate
\[
 \boxed{
 \beta_{\op}
 \le
 \min\left\{
 16,
 16P_d
 +\frac{2^{\alpha+4}(L_N+10)}{d}
 +\frac{32T_d}{d^2}\,\Xi_{\alpha,J}P_d
 \right\}
 }
\]
also holds.  Here \(N>d\) implies \(m=d\), so \(P_d=P_m\).
\end{theorem}

\begin{remark}[When the probability bound is informative]
The theorem is a finite-dimensional statement and does not require a separate
relation between \(d\) and \(N\).  If the displayed lower bound on probability
is negative, it is simply vacuous.  In the common regime \(N=d^{O(1)}\), every
term in \eqref{eq:total-failure} is \(e^{-\Omega(d)}\), because
\(K\le d/256\) and \(J=O(\log d)\).
\end{remark}

\section{Exact Hessian formula and norm geometry}
\label{sec:hessian}

For \(i\in[N]\), write
\[
 q_i(W):=(q_{i1}(W),\ldots,q_{iN}(W))^\top\in\R^N
\]
and define
\[
 C_i(W)
 :=
 U\bigl(\Diag(q_i(W))-q_i(W)q_i(W)^\top\bigr)U^\top.
\]

\begin{proposition}[Gradient and Hessian]
\label{prop:hessian}
For every \(W,A,B\in\R^{d\times d}\),
\[
 \nabla \ell_i(W)
 =
 \left(\sum_{j=1}^N q_{ij}(W)u_j-u_i\right)v_i^\top,
\]
\[
 \ip{B}{\nabla^2L(W)[A]}
 =
 \sum_{i=1}^N p_i(Bv_i)^\top C_i(W)(Av_i).
\]
Moreover, \(C_i(W)\succeq 0\), and hence
\[
 \boxed{
 \nabla^2L(W)[H,H]
 =
 \sum_{i=1}^N p_i(Hv_i)^\top C_i(W)(Hv_i)
 \ge 0.
 }
\]
\end{proposition}

\begin{proof}
For a direction \(A\),
\[
 Dz_{ij}(W)[A]=u_j^\top Av_i.
\]
Therefore
\begin{align*}
 D\ell_i(W)[A]
 &=
 \sum_{j=1}^N q_{ij}(W)u_j^\top Av_i-u_i^\top Av_i\\
 &=
 \left(\sum_{j=1}^N q_{ij}(W)u_j-u_i\right)^\top Av_i.
\end{align*}
Using the Frobenius inner product gives the displayed gradient formula.

Let
\[
 \mu_i(W):=\sum_{j=1}^N q_{ij}(W)u_j=Uq_i(W).
\]
The Jacobian of the softmax map is
\[
 D\softmax(z)=\Diag(q)-qq^\top.
\]
Consequently,
\[
 Dq_i(W)[A]
 =
 \bigl(\Diag(q_i(W))-q_i(W)q_i(W)^\top\bigr)U^\top Av_i,
\]
and hence
\[
 D\mu_i(W)[A]=C_i(W)Av_i.
\]
Because the target term is linear,
\[
 \nabla^2\ell_i(W)[A]=C_i(W)Av_iv_i^\top.
\]
Summing over \(i\) and taking the Frobenius inner product with \(B\) yields
\[
 \ip{B}{\nabla^2L(W)[A]}
 =
 \sum_{i=1}^N p_i(Bv_i)^\top C_i(W)(Av_i).
\]
Finally,
\[
 C_i(W)
 =
 \sum_{j=1}^N q_{ij}(W)u_ju_j^\top
 -\left(\sum_{j=1}^N q_{ij}(W)u_j\right)
  \left(\sum_{j=1}^N q_{ij}(W)u_j\right)^\top
\]
is a covariance matrix, so it is positive semidefinite.
\end{proof}

\begin{proposition}[Quadratic characterization of smoothness]
\label{prop:quadratic-characterization}
Let \(\norm{\cdot}_X\) be any norm on \(\R^{d\times d}\), and let
\(\norm{\cdot}_{X,*}\) be its dual norm.  For every fixed \(W\),
\[
 \sup_{\norm{H}_X\le 1}
 \norm{\nabla^2L(W)[H]}_{X,*}
 =
 \sup_{\norm{H}_X\le 1}
 \nabla^2L(W)[H,H].
\]
Consequently,
\[
 \sup_W\sup_{\norm{H}_X\le 1}
 \nabla^2L(W)[H,H]
\]
is the smallest number \(\beta\) such that
\[
 \norm{\nabla L(W_1)-\nabla L(W_0)}_{X,*}
 \le
 \beta\norm{W_1-W_0}_X
\]
for all \(W_0,W_1\in\R^{d\times d}\).
\end{proposition}

\begin{proof}
For fixed \(W\), define the symmetric positive-semidefinite bilinear form
\[
 B_W(A,H):=\ip{A}{\nabla^2L(W)[H]}.
\]
Positive semidefiniteness follows from Proposition~\ref{prop:hessian}.  Applying
nonnegativity to \(B_W(A+tH,A+tH)\) as a quadratic polynomial in \(t\) gives
\[
 \abs{B_W(A,H)}^2\le B_W(A,A)B_W(H,H).
\]
Therefore
\begin{align*}
 \sup_{\norm{H}_X\le1}
 \norm{\nabla^2L(W)[H]}_{X,*}
 &=
 \sup_{\norm{A}_X\le1,\ \norm{H}_X\le1}
 \abs{B_W(A,H)}\\
 &\le
 \sup_{\norm{G}_X\le1}B_W(G,G).
\end{align*}
The reverse inequality follows by choosing \(A=H=G\).

For the Lipschitz statement, use the fundamental theorem of calculus:
\[
 \nabla L(W_1)-\nabla L(W_0)
 =
 \int_0^1
 \nabla^2L\bigl(W_0+t(W_1-W_0)\bigr)[W_1-W_0] \,dt.
\]
The preceding operator-norm identity gives the upper bound.  Conversely, a
global gradient Lipschitz bound implies the same Hessian bound by taking a
directional derivative at any \(W\).
\end{proof}

\section{Deterministic reductions}
\label{sec:deterministic}

Define
\[
 R_u:=\max_{j\in[N]}\norm{u_j}^2,
 \qquad
 M_v:=\sum_{i=1}^N p_i v_iv_i^\top.
\]
For \(y\in\R^d\), define the output width
\[
 w_U(y)
 :=
 \max_{j,k\in[N]}(u_j-u_k)^\top y
 =
 \max_j u_j^\top y-\min_k u_k^\top y.
\]

\begin{lemma}[Uniform covariance bound]
\label{lem:covariance-upper}
For every probability vector \(q\in\R^N\), define
\[
 C(q)
 :=
 \sum_{j=1}^N q_j u_ju_j^\top
 -\left(\sum_{j=1}^N q_j u_j\right)
  \left(\sum_{j=1}^N q_j u_j\right)^\top.
\]
Then
\[
 0\preceq C(q)\preceq R_u I_d.
\]
\end{lemma}

\begin{proof}
For every \(x\in\R^d\),
\begin{align*}
 x^\top C(q)x
 &=
 \sum_j q_j(x^\top u_j)^2
 -\left(\sum_j q_jx^\top u_j\right)^2\\
 &\le
 \sum_j q_j(x^\top u_j)^2
 \le
 R_u\norm{x}^2.
\end{align*}
The lower bound follows because \(C(q)\) is a covariance matrix.
\end{proof}

\begin{lemma}[Frobenius reduction]
\label{lem:frobenius-reduction}
For every realization of the embeddings,
\[
 \boxed{\beta_F\le R_u\norm{M_v}_{\op}.}
\]
\end{lemma}

\begin{proof}
By Proposition~\ref{prop:hessian} and
Lemma~\ref{lem:covariance-upper},
\begin{align*}
 \nabla^2L(W)[H,H]
 &\le
 R_u\sum_{i=1}^N p_i\norm{Hv_i}^2\\
 &=
 R_u\Tr(H^\top HM_v)\\
 &\le
 R_u\norm{M_v}_{\op}\norm{H}_F^2.
\end{align*}
Take the supremum over \(W\) and nonzero \(H\).
\end{proof}

\begin{lemma}[Operator-width reduction]
\label{lem:operator-reduction}
For every realization,
\[
 \boxed{
 \beta_{\op}
 \le
 \frac14
 \sup_{\norm{H}_{\op}\le1}
 \sum_{i=1}^N p_i w_U(Hv_i)^2.
 }
\]
Also,
\[
 \boxed{\beta_{\op}\le R_u\Tr(M_v).}
\]
\end{lemma}

\begin{proof}
Let a scalar random variable \(X\) be supported in \([a,b]\).  Then
\[
 \Var(X)
 =\min_c\E(X-c)^2
 \le
 \E\left(X-\frac{a+b}{2}\right)^2
 \le
 \frac{(b-a)^2}{4}.
\]
For fixed \(y\), apply this to the random variable \(u_j^\top y\) under
\(j\sim q_i(W)\).  Its range length is \(w_U(y)\), so
\[
 y^\top C_i(W)y\le\frac14w_U(y)^2.
\]
Substitute \(y=Hv_i\) into Proposition~\ref{prop:hessian} and then take the
supremum over \(W\) and \(\norm{H}_{\op}\le1\).

The second estimate follows from
\begin{align*}
 \nabla^2L(W)[H,H]
 &\le
 R_u\sum_i p_i\norm{Hv_i}^2\\
 &\le
 R_u\norm{H}_{\op}^2\sum_i p_i\norm{v_i}^2
 =
 R_u\norm{H}_{\op}^2\Tr(M_v).
\end{align*}
\end{proof}

\section{Probability tools}
\label{sec:probability-tools}

\begin{lemma}[Weighted chi-square inequality]
\label{lem:weighted-chi-square}
Let \(\xi_1,\ldots,\xi_M\) be independent standard normal variables, and let
\(a_1,\ldots,a_M\ge0\).  Define
\[
 S:=\sum_{r=1}^M a_r(\xi_r^2-1).
\]
Then, for every \(t>0\),
\[
 \Pp\left(
 S\ge2\norm{a}_2\sqrt t+2\norm{a}_\infty t
 \right)\le e^{-t},
\]
\[
 \Pp\left(
 S\le-2\norm{a}_2\sqrt t
 \right)\le e^{-t}.
\]
\end{lemma}

\begin{proof}
If \(\norm{a}_2=0\), then every \(a_r=0\), so \(S=0\) and both claims are
immediate.  Assume \(\norm{a}_2>0\).

For \(0<\lambda<(2\norm{a}_\infty)^{-1}\),
\[
 \E e^{\lambda a_r(\xi_r^2-1)}
 =e^{-\lambda a_r}(1-2\lambda a_r)^{-1/2}.
\]
For \(0\le x<1/2\),
\[
 -\frac12\log(1-2x)-x\le\frac{x^2}{1-2x}.
\]
By independence,
\[
 \log\E e^{\lambda S}
 \le
 \frac{\lambda^2\norm{a}_2^2}
 {1-2\lambda\norm{a}_\infty}.
\]
Chernoff's inequality with
\[
 \lambda
 =
 \frac{\sqrt t}{\norm{a}_2+2\norm{a}_\infty\sqrt t}
\]
gives the upper tail.

For the lower tail,
\[
 \E e^{-\lambda a_r(\xi_r^2-1)}
 =e^{\lambda a_r}(1+2\lambda a_r)^{-1/2}.
\]
Since
\[
 x-\frac12\log(1+2x)\le x^2
 \qquad (x\ge0),
\]
we have
\[
 \log\E e^{-\lambda S}\le\lambda^2\norm{a}_2^2.
\]
Chernoff's inequality with \(\lambda=\sqrt t/\norm{a}_2\) gives the lower
tail.
\end{proof}

\begin{lemma}[Covering nets]
\label{lem:nets}
The Euclidean unit sphere in \(\R^q\) has an \(r\)-net in Euclidean norm of
cardinality at most
\[
 \left(1+\frac2r\right)^q.
\]
The operator-norm unit ball in \(\R^{d\times d}\) has a
\(d^{-1/2}\)-net in operator norm of cardinality at most
\[
 (3d)^{d^2}.
\]
\end{lemma}

\begin{proof}
A maximal \(r\)-separated subset of a normed set is an \(r\)-net.  For the
Euclidean sphere, the Euclidean balls of radius \(r/2\) centered at the net
points are disjoint and contained in the Euclidean ball of radius \(1+r/2\).
Volume comparison gives the first claim.

For the second claim, take a maximal \(d^{-1/2}\)-separated subset in operator
norm of the operator unit ball.  Since
\[
 \norm{A}_F\le\sqrt d\norm{A}_{\op},
\]
the operator unit ball lies in the Frobenius ball of radius \(\sqrt d\) in
ambient dimension \(d^2\).  Moreover, operator-norm separation implies
Frobenius-norm separation.  Euclidean volume comparison therefore gives
\[
 \left(1+\frac{2\sqrt d}{d^{-1/2}}\right)^{d^2}
 =(1+2d)^{d^2}
 \le(3d)^{d^2}.
\]
\end{proof}

\begin{lemma}[Gaussian vector estimates]
\label{lem:gaussian-vector}
Let \(g,h\) be independent \(\mathcal N(0,I_d/d)\) vectors.  Then
\[
 \Pp\left(\norm{g}^2\notin[1/2,3/2]\right)
 \le2e^{-d/32},
\]
\[
 \Pp(\norm{g}>2)\le e^{-d/2},
\]
and
\[
 \Pp\left(
 \abs{g^\top h}>\frac14,
 \ \norm{g}^2\le\frac32
 \right)
 \le2e^{-d/48}.
\]
\end{lemma}

\begin{proof}
Because \(d\norm{g}^2\sim\chi_d^2\), Chernoff optimization gives
\[
 \Pp(\chi_d^2\ge ad)
 \le
 \exp\left[-\frac d2(a-1-\log a)\right]
 \qquad (a>1),
\]
and
\[
 \Pp(\chi_d^2\le ad)
 \le
 \exp\left[-\frac d2(a-1-\log a)\right]
 \qquad (0<a<1).
\]
For \(a=3/2\) and \(a=1/2\), the exponent is at least \(d/32\).  For
\(a=4\), the exponent is at least \(d/2\).  This proves the first two claims.

Conditional on \(g\),
\[
 g^\top h\sim\mathcal N\left(0,\frac{\norm{g}^2}{d}\right).
\]
On the event \(\norm{g}^2\le3/2\), the two-sided Gaussian tail bound gives
\[
 \Pp\left(\abs{g^\top h}>\frac14\mid g\right)
 \le
 2\exp\left(-\frac{(1/4)^2}{2(3/2)/d}\right)
 =2e^{-d/48}.
\]
\end{proof}

\begin{lemma}[Rectangular Gaussian conditioning]
\label{lem:rectangular-gaussian}
Let \(G\in\R^{d\times K}\) have independent \(\mathcal N(0,1/d)\) entries,
and assume \(1\le K\le d/256\).  Then
\[
 \Pp\left(
 \sigma_{\min}(G)<\frac12
 \ \text{or}\ 
 \norm{G}_{\op}>\frac32
 \right)
 \le
 2\cdot17^K e^{-d/32}.
\]
\end{lemma}

\begin{proof}
Take a \(1/8\)-net \(\mathcal N\) of \(\sphere^{K-1}\) with
\(\abs{\mathcal N}\le17^K\).  For every fixed \(x\in\sphere^{K-1}\),
\[
 Gx\sim\mathcal N(0,I_d/d).
\]
Lemma~\ref{lem:gaussian-vector} and a union bound imply that, except on an
event of probability at most \(2\cdot17^K e^{-d/32}\),
\[
 \frac1{\sqrt2}\le\norm{Gx_0}\le\sqrt{\frac32}
 \qquad\text{for every }x_0\in\mathcal N.
\]
Assume this event.  For arbitrary \(x\in\sphere^{K-1}\), choose
\(x_0\in\mathcal N\) with \(\norm{x-x_0}\le1/8\).  First,
\[
 \norm{G}_{\op}
 \le
 \sqrt{\frac32}+\frac18\norm{G}_{\op},
\]
so
\[
 \norm{G}_{\op}
 \le
 \frac87\sqrt{\frac32}<\frac32.
\]
Second,
\[
 \norm{Gx}
 \ge
 \frac1{\sqrt2}
 -\frac18\norm{G}_{\op}
 \ge
 \frac1{\sqrt2}
 -\frac17\sqrt{\frac32}
 >\frac12.
\]
Taking the infimum over \(x\in\sphere^{K-1}\) proves the lower singular-value
bound.
\end{proof}

\section{Frobenius upper bound}
\label{sec:frobenius-upper}

\begin{lemma}[Uniform upper bound for the weighted input covariance]
\label{lem:Mv-upper}
Define the event
\[
 \mathcal E_M
 :=
 \left\{
 \norm{M_v}_{\op}\le\frac6d+20p_1
 \right\}.
\]
Then
\[
 \Pp(\mathcal E_M)\ge1-e^{-d}.
\]
\end{lemma}

\begin{proof}
Fix \(x\in\sphere^{d-1}\).  Since
\(\sqrt d\,x^\top v_i\) are independent standard normal variables, there are
independent \(\xi_i\sim\mathcal N(0,1)\) such that
\[
 x^\top M_vx
 =
 \frac1d\sum_{i=1}^N p_i\xi_i^2.
\]
Lemma~\ref{lem:weighted-chi-square} gives, for every \(t>0\),
\[
 x^\top M_vx
 \le
 \frac1d\left(1+2\norm{p}_2\sqrt t+2p_1t\right)
\]
with failure probability at most \(e^{-t}\).

Take a \(1/4\)-net \(\mathcal N\) of \(\sphere^{d-1}\) with
\(\abs{\mathcal N}\le9^d\), and choose
\[
 t=d(\log9+1).
\]
A union bound shows that the preceding estimate holds for every
\(x\in\mathcal N\) with failure probability at most
\[
 9^d e^{-d(\log9+1)}=e^{-d}.
\]
Assume this net event.  Let \(y\) be a unit top eigenvector of \(M_v\), and
choose \(x\in\mathcal N\) with \(\norm{x-y}\le1/4\).  Since \(M_v\succeq0\),
\begin{align*}
 \norm{M_v}_{\op}
 &=y^\top M_vy\\
 &\le
 x^\top M_vx
 +\abs{(y-x)^\top M_vy}
 +\abs{x^\top M_v(y-x)}\\
 &\le
 x^\top M_vx+\frac12\norm{M_v}_{\op}.
\end{align*}
Thus
\[
 \norm{M_v}_{\op}
 \le
 \frac2d\left(1+2\norm{p}_2\sqrt t+2p_1t\right).
\]
Because
\[
 \norm{p}_2^2=\sum_i p_i^2\le p_1\sum_i p_i=p_1
\]
and \(t<4d\),
\[
 \norm{M_v}_{\op}
 \le
 \frac2d+8\sqrt{\frac{p_1}{d}}+16p_1.
\]
Using \(8\sqrt{ab}\le4a+4b\) gives
\[
 \norm{M_v}_{\op}\le\frac6d+20p_1.
\]
\end{proof}

\begin{proposition}[Frobenius upper estimate on explicit events]
\label{prop:frobenius-upper}
On the event
\[
 \left\{\max_{j\in[N]}\norm{u_j}\le2\right\}\cap\mathcal E_M,
\]
one has
\[
 \beta_F\le80\left(\frac1d+p_1\right).
\]
\end{proposition}

\begin{proof}
On the displayed event, \(R_u\le4\).  Lemmas
\ref{lem:frobenius-reduction} and \ref{lem:Mv-upper} give
\[
 \beta_F
 \le
 4\left(\frac6d+20p_1\right)
 =
 \frac{24}{d}+80p_1
 \le
 80\left(\frac1d+p_1\right).
\]
\end{proof}

\section{Two-output selectors}
\label{sec:selectors}

\begin{lemma}[Deterministic two-output selector]
\label{lem:deterministic-selector}
Suppose two output vectors \(u_a,u_b\) satisfy
\[
 \frac12\le\norm{u_a}^2,\norm{u_b}^2\le\frac32,
 \qquad
 \abs{u_a^\top u_b}\le\frac14.
\]
Then there exists \(y\in\R^d\) such that
\[
 u_a^\top y=u_b^\top y=1,
 \qquad
 \norm{y}^2\le8.
\]
If, in addition,
\[
 u_j^\top y\le\frac12
 \qquad\text{for every }j\notin\{a,b\},
\]
then, as \(t\to\infty\),
\[
 \softmax(tU^\top y)
 \longrightarrow
 \frac12e_a+\frac12e_b,
\]
and the associated output covariance converges to
\[
 C_\infty
 =
 \frac14(u_a-u_b)(u_a-u_b)^\top.
\]
For the unit vector
\[
 x:=\frac{u_a-u_b}{\norm{u_a-u_b}},
\]
one has
\[
 x^\top C_\infty x\ge\frac18.
\]
\end{lemma}

\begin{proof}
Let
\[
 A:=[u_a,u_b],
 \qquad
 G:=A^\top A.
\]
By the Gershgorin bound for the symmetric \(2\times2\) matrix \(G\),
\[
 \lambda_{\min}(G)
 \ge
 \min\{\norm{u_a}^2,\norm{u_b}^2\}
 -\abs{u_a^\top u_b}
 \ge\frac14.
\]
Thus \(G\) is invertible.  Define
\[
 y:=AG^{-1}\binom11.
\]
Then
\[
 A^\top y=GG^{-1}\binom11=\binom11,
\]
and
\[
 \norm{y}^2
 =(1,1)G^{-1}(1,1)^\top
 \le
 \frac{\norm{(1,1)}^2}{\lambda_{\min}(G)}
 \le8.
\]
The selected logits are both equal to \(t\), while every other logit is at
most \(t/2\).  Dividing numerator and denominator of the softmax by \(e^t\)
proves the softmax limit.  The covariance of the limiting two-point uniform
distribution is
\begin{align*}
 C_\infty
 &=
 \frac12u_au_a^\top+\frac12u_bu_b^\top
 -\frac14(u_a+u_b)(u_a+u_b)^\top\\
 &=
 \frac14(u_a-u_b)(u_a-u_b)^\top.
\end{align*}
Finally,
\[
 \norm{u_a-u_b}^2
 =
 \norm{u_a}^2+\norm{u_b}^2-2u_a^\top u_b
 \ge
 \frac12+\frac12-2\cdot\frac14
 =\frac12,
\]
so
\[
 x^\top C_\infty x
 =\frac14\norm{u_a-u_b}^2
 \ge\frac18.
\]
\end{proof}

\begin{lemma}[Simultaneous selector event]
\label{lem:simultaneous-selector}
Fix \(K\ge1\) disjoint output pairs
\[
 (u_1,u_2),\ldots,(u_{2K-1},u_{2K}),
 \qquad 2K\le N.
\]
There is an event \(\mathcal E_{\mathrm{sel}}\) such that
\[
 \Pp(\mathcal E_{\mathrm{sel}})
 \ge
 1-
 \left(
 4K e^{-d/32}
 +2K e^{-d/48}
 +K(N-2)e^{-d/64}
 \right),
\]
and, on \(\mathcal E_{\mathrm{sel}}\), every pair admits vectors
\(y_r,x_r\) and a covariance limit \(C_{\infty,r}\) satisfying all
conclusions of Lemma~\ref{lem:deterministic-selector}.
\end{lemma}

\begin{proof}
For each of the \(2K\) paired vectors, the norm condition
\(\norm{u_j}^2\in[1/2,3/2]\) fails with probability at most
\(2e^{-d/32}\).  Hence all paired norm conditions hold except on an event of
probability at most
\[
 4K e^{-d/32}.
\]

For each pair \((u_{2r-1},u_{2r})\), Lemma~\ref{lem:gaussian-vector} gives
\[
 \Pp\left(
 \abs{u_{2r-1}^\top u_{2r}}>\frac14,
 \ \norm{u_{2r-1}}^2\le\frac32
 \right)
 \le2e^{-d/48}.
\]
Thus the probability of a pairwise inner-product failure while the relevant
norm upper bound holds is at most \(2K e^{-d/48}\).

On the complement of these two failure events,
Lemma~\ref{lem:deterministic-selector} constructs \(y_r\), depending only on
\((u_{2r-1},u_{2r})\), with \(\norm{y_r}^2\le8\).  For
\(j\notin\{2r-1,2r\}\), the vector \(u_j\) is independent of this pair, and
therefore, conditionally on the pair,
\[
 u_j^\top y_r
 \sim
 \mathcal N\left(0,\frac{\norm{y_r}^2}{d}\right).
\]
The one-sided Gaussian tail bound implies
\[
 \Pp\left(u_j^\top y_r>\frac12\mid u_{2r-1},u_{2r}\right)
 \le
 \exp\left(-\frac{(1/2)^2}{2(8/d)}\right)
 =e^{-d/64}.
\]
Taking conditional expectations and union bounding over the \(K(N-2)\)
pair--outside-index choices gives the third failure term.  On the resulting
event, every pair satisfies the hypotheses of
Lemma~\ref{lem:deterministic-selector}.
\end{proof}

\section{Frobenius lower bound}
\label{sec:frobenius-lower}

Let
\[
 V_K:=[v_1,\ldots,v_K]\in\R^{d\times K},
\]
and define the rectangular input event
\[
 \mathcal E_V
 :=
 \left\{
 \sigma_{\min}(V_K)\ge\frac12,
 \ \norm{V_K}_{\op}\le\frac32
 \right\}.
\]
By Lemma~\ref{lem:rectangular-gaussian},
\begin{equation}
 \Pp(\mathcal E_V)
 \ge1-2\cdot17^K e^{-d/32}.
 \label{eq:EV-probability}
\end{equation}

\begin{lemma}[Single-atom Frobenius lower bound]
\label{lem:single-atom-lower}
On \(\mathcal E_{\mathrm{sel}}\cap\mathcal E_V\),
\[
 \beta_F\ge\frac{p_1}{32}.
\]
\end{lemma}

\begin{proof}
Use the selector \(y_1,x_1,C_{\infty,1}\) for the first output pair
\((u_1,u_2)\).  On \(\mathcal E_V\),
\[
 \norm{v_1}=\norm{V_Ke_1}\ge\sigma_{\min}(V_K)\ge\frac12,
\]
so \(\norm{v_1}^2\ge1/4\).  Define
\[
 W_t:=t\,y_1\frac{v_1^\top}{\norm{v_1}^2},
 \qquad
 H:=x_1\frac{v_1^\top}{\norm{v_1}}.
\]
Then
\[
 W_tv_1=ty_1,
 \qquad
 \norm{H}_F=1,
 \qquad
 Hv_1=x_1\norm{v_1}.
\]
All Hessian summands are nonnegative.  Therefore
\[
 \nabla^2L(W_t)[H,H]
 \ge
 p_1\norm{v_1}^2x_1^\top C_1(W_t)x_1.
\]
As \(t\to\infty\), the selector property gives
\(C_1(W_t)\to C_{\infty,1}\).  Hence
\[
 \beta_F
 \ge
 p_1\norm{v_1}^2x_1^\top C_{\infty,1}x_1
 \ge
 p_1\cdot\frac14\cdot\frac18
 =\frac{p_1}{32}.
\]
\end{proof}

\begin{lemma}[Weighted trace lower bound]
\label{lem:trace-lower}
Define
\[
 \mathcal E_{\mathrm{tr}}
 :=
 \left\{\Tr(M_v)\ge\frac12\right\}.
\]
Then
\[
 \Pp(\mathcal E_{\mathrm{tr}})\ge1-e^{-d/16}.
\]
\end{lemma}

\begin{proof}
Write
\[
 \Tr(M_v)
 =
 \frac1d\sum_{i=1}^N\sum_{k=1}^d p_i g_{ik}^2,
\]
where all \(g_{ik}\) are independent standard normal variables.  Apply the
lower-tail part of Lemma~\ref{lem:weighted-chi-square} to the weights
\[
 a_{ik}:=\frac{p_i}{d}.
\]
Their Euclidean norm is
\[
 \norm{a}_2
 =
 \frac{\norm{p}_2}{\sqrt d}.
\]
Take
\[
 t:=\frac{d}{16\norm{p}_2^2}.
\]
Then \(2\norm{a}_2\sqrt t=1/2\), and therefore
\[
 \Pp\left(\Tr(M_v)<\frac12\right)
 \le
 \exp\left(-\frac{d}{16\norm{p}_2^2}\right)
 \le e^{-d/16},
\]
because \(\norm{p}_2^2\le(\sum_i p_i)^2=1\).
\end{proof}

\begin{lemma}[Diffuse Frobenius lower bound]
\label{lem:diffuse-lower}
On \(\mathcal E_{\mathrm{sel}}\cap\mathcal E_{\mathrm{tr}}\),
\[
 \beta_F\ge\frac1{32d}.
\]
\end{lemma}

\begin{proof}
Use the selector \(y_1,x_1,C_{\infty,1}\) for the first output pair.  Gaussian
input vectors are nonzero almost surely; we work on this full-probability
event.  For \(b\in\sphere^{d-1}\), define
\[
 M_+(b)
 :=
 \sum_{i:\,b^\top v_i>0}p_i v_iv_i^\top.
\]
Let \(b\) be uniform on \(\sphere^{d-1}\), independently of the embeddings.
For every fixed nonzero \(v_i\), symmetry gives
\[
 \Pp_b(b^\top v_i>0)=\frac12.
\]
Thus
\[
 \E_b\Tr(M_+(b))
 =
 \frac12\Tr(M_v).
\]
The finite union of hyperplanes \(\{b:b^\top v_i=0\}\) has spherical measure
zero.  Hence there exists a \(b\) outside all these hyperplanes such that
\[
 \Tr(M_+(b))\ge\frac12\Tr(M_v).
\]
On \(\mathcal E_{\mathrm{tr}}\), this gives
\[
 \Tr(M_+(b))\ge\frac14.
\]
Let \(a\in\sphere^{d-1}\) be a top eigenvector of \(M_+(b)\).  Since a
positive-semidefinite \(d\times d\) matrix has largest eigenvalue at least its
trace divided by \(d\),
\[
 a^\top M_+(b)a\ge\frac1{4d}.
\]
Define
\[
 W_t:=t y_1b^\top,
 \qquad
 H:=x_1a^\top.
\]
Then \(\norm{H}_F=1\).  If \(b^\top v_i>0\), then
\[
 W_tv_i=t(b^\top v_i)y_1,
\]
and the scalar \(t(b^\top v_i)\) tends to \(+\infty\).  Hence
\[
 C_i(W_t)\longrightarrow C_{\infty,1}.
\]
Using nonnegativity of all Hessian terms and the finiteness of the selected
index set,
\begin{align*}
 \beta_F
 &\ge
 \liminf_{t\to\infty}\nabla^2L(W_t)[H,H]\\
 &\ge
 \sum_{i:\,b^\top v_i>0}
 p_i(a^\top v_i)^2x_1^\top C_{\infty,1}x_1\\
 &=
 x_1^\top C_{\infty,1}x_1\,a^\top M_+(b)a\\
 &\ge
 \frac18\cdot\frac1{4d}
 =\frac1{32d}.
\end{align*}
\end{proof}

\begin{proposition}[Frobenius lower estimate on explicit events]
\label{prop:frobenius-lower}
On
\[
 \mathcal E_{\mathrm{sel}}
 \cap\mathcal E_V
 \cap\mathcal E_{\mathrm{tr}},
\]
one has
\[
 \beta_F
 \ge
 \frac1{64}\left(p_1+\frac1d\right).
\]
\end{proposition}

\begin{proof}
Lemmas~\ref{lem:single-atom-lower} and \ref{lem:diffuse-lower} give
\[
 \beta_F
 \ge
 \max\left\{\frac{p_1}{32},\frac1{32d}\right\}.
\]
Since \(\max\{a,b\}\ge(a+b)/2\),
\[
 \beta_F
 \ge
 \frac1{64}\left(p_1+\frac1d\right).
\]
\end{proof}

\section{Operator-norm lower bound}
\label{sec:operator-lower}

Define, for \(i\in[K]\),
\[
 g_i:=\frac{u_{2i-1}-u_{2i}}{\sqrt2},
 \qquad
 G_K:=[g_1,\ldots,g_K].
\]
The columns \(g_i\) are independent \(\mathcal N(0,I_d/d)\) vectors.  Define
\[
 \mathcal E_G
 :=
 \left\{
 \sigma_{\min}(G_K)\ge\frac12,
 \ \norm{G_K}_{\op}\le\frac32
 \right\}.
\]
Lemma~\ref{lem:rectangular-gaussian} gives
\begin{equation}
 \Pp(\mathcal E_G)
 \ge1-2\cdot17^K e^{-d/32}.
 \label{eq:EG-probability}
\end{equation}

\begin{proposition}[Simultaneous activation]
\label{prop:simultaneous-activation}
On
\[
 \mathcal E_{\mathrm{sel}}\cap\mathcal E_V\cap\mathcal E_G,
\]
one has
\[
 \beta_{\op}\ge\frac1{288}\sum_{i=1}^K p_i.
\]
\end{proposition}

\begin{proof}
For each \(i\in[K]\), let \(y_i,x_i,C_{\infty,i}\) be the selector objects
for the output pair \((u_{2i-1},u_{2i})\).  By construction,
\[
 x_i=\frac{u_{2i-1}-u_{2i}}{\norm{u_{2i-1}-u_{2i}}}
 =\frac{g_i}{\norm{g_i}},
\]
and
\[
 x_i^\top C_{\infty,i}x_i\ge\frac18.
\]
Let
\[
 X_K:=[x_1,\ldots,x_K],
 \qquad
 Y_K:=[y_1,\ldots,y_K].
\]
On \(\mathcal E_G\), every \(\norm{g_i}\ge\sigma_{\min}(G_K)\ge1/2\).
Writing
\[
 D:=\Diag(\norm{g_1},\ldots,\norm{g_K}),
 \qquad
 X_K=G_KD^{-1},
\]
we obtain
\[
 \norm{X_K}_{\op}
 \le
 \norm{G_K}_{\op}\norm{D^{-1}}_{\op}
 \le
 \frac32\cdot2=3.
\]

On \(\mathcal E_V\), \(V_K\) has full column rank and
\[
 \norm{V_K^\dagger}_{\op}
 =\frac1{\sigma_{\min}(V_K)}
 \le2,
\]
where \(V_K^\dagger=(V_K^\top V_K)^{-1}V_K^\top\).  Define
\[
 W_t:=tY_KV_K^\dagger.
\]
Since \(V_K^\dagger V_K=I_K\),
\[
 W_tv_i=ty_i
 \qquad (i\in[K]).
\]
Therefore
\[
 C_i(W_t)\longrightarrow C_{\infty,i}
 \qquad (i\in[K]).
\]
Now define
\[
 H:=\frac16X_KV_K^\dagger.
\]
Then
\[
 \norm{H}_{\op}
 \le
 \frac16\norm{X_K}_{\op}\norm{V_K^\dagger}_{\op}
 \le1,
\]
and
\[
 Hv_i=\frac16x_i
 \qquad (i\in[K]).
\]
Using nonnegativity of the Hessian summands,
\begin{align*}
 \beta_{\op}
 &\ge
 \liminf_{t\to\infty}\nabla^2L(W_t)[H,H]\\
 &\ge
 \sum_{i=1}^K p_i
 \left(\frac16x_i\right)^\top
 C_{\infty,i}
 \left(\frac16x_i\right)\\
 &\ge
 \frac1{36}\cdot\frac18\sum_{i=1}^Kp_i
 =
 \frac1{288}\sum_{i=1}^Kp_i.
\end{align*}
\end{proof}

\begin{lemma}[Head-mass comparison]
\label{lem:head-mass}
Let \(a_1\ge\cdots\ge a_m\ge0\), and let \(1\le K\le m\).  Then
\[
 \sum_{i=1}^K a_i
 \ge
 \frac Km\sum_{i=1}^m a_i.
\]
\end{lemma}

\begin{proof}
If \(K=m\), the claim is equality.  Assume \(K<m\).  Since
\[
 a_{K+1}\le a_K\le\frac1K\sum_{i=1}^K a_i,
\]
we have
\[
 \sum_{i=K+1}^m a_i
 \le
 (m-K)a_{K+1}
 \le
 \frac{m-K}{K}\sum_{i=1}^K a_i.
\]
Adding the head mass to both sides and rearranging proves the result.
\end{proof}

\begin{proposition}[Operator lower estimate on explicit events]
\label{prop:operator-lower}
Assume \(m\ge512\).  On
\[
 \mathcal E_{\mathrm{sel}}
 \cap\mathcal E_V
 \cap\mathcal E_G
 \cap\mathcal E_{\mathrm{tr}},
\]
one has
\[
 \beta_{\op}
 \ge
 \frac1{294912}\left(P_m+\frac1d\right).
\]
\end{proposition}

\begin{proof}
Because \(m/256\ge2\),
\[
 K=\left\lfloor\frac{m}{256}\right\rfloor
 \ge\frac{m}{512}.
\]
Apply Lemma~\ref{lem:head-mass} to the nonincreasing sequence
\(p_1,\ldots,p_m\).  Proposition~\ref{prop:simultaneous-activation} gives
\[
 \beta_{\op}
 \ge
 \frac1{288}\sum_{i=1}^Kp_i
 \ge
 \frac1{288}\frac KmP_m
 \ge
 \frac{P_m}{147456}.
\]
Moreover, the Frobenius unit ball is contained in the operator-norm unit ball,
because \(\norm{H}_{\op}\le\norm{H}_F\).  Hence
\[
 \beta_{\op}\ge\beta_F.
\]
On the displayed event, Lemma~\ref{lem:diffuse-lower} therefore gives
\[
 \beta_{\op}\ge\frac1{32d}.
\]
Combining the two estimates,
\begin{align*}
 \beta_{\op}
 &\ge
 \max\left\{\frac{P_m}{147456},\frac1{32d}\right\}\\
 &\ge
 \frac{P_m}{294912}+\frac1{64d}\\
 &\ge
 \frac1{294912}\left(P_m+\frac1d\right).
\end{align*}
\end{proof}

\section{Uniform spectral-width estimate}
\label{sec:block-width}

\begin{lemma}[Uniform block spectral-width bound]
\label{lem:block-width}
Let \(I\subseteq[N]\) be deterministic, with \(s:=\abs{I}\ge1\).  Fix an
output realization satisfying
\[
 \max_{j\in[N]}\norm{u_j}\le2.
\]
For \(\tau>0\), define
\[
 T(\tau):=d^2\log(3d)+\tau,
 \qquad
 b_s(\tau):=\max\left\{8,\frac{2T(\tau)}s\right\}.
\]
Conditionally on this output realization, there is an event
\(\mathcal E_I(\tau)\), depending only on \(\{v_i:i\in I\}\), such that
\[
 \Pp\bigl(\mathcal E_I(\tau)\mid U\bigr)\ge1-e^{-\tau},
\]
and, on
\[
 \mathcal E_I(\tau)
 \cap
 \left\{\max_{i\in I}\norm{v_i}\le2\right\},
\]
one has
\[
 \boxed{
 \sup_{\norm{H}_{\op}\le1}
 \frac1s\sum_{i\in I}w_U(Hv_i)^2
 \le
 \frac{64}{d}\left(L_N+b_s(\tau)+2\right).
 }
\]
\end{lemma}

\begin{proof}
Fix \(H\) with \(\norm{H}_{\op}\le1\), and let
\(v\sim\mathcal N(0,I_d/d)\), independently of \(U\).  For fixed \(j,k\),
conditionally on \(U\),
\[
 (u_j-u_k)^\top Hv
 \sim
 \mathcal N\left(
 0,\frac{\norm{H^\top(u_j-u_k)}^2}{d}
 \right).
\]
Because \(\norm{u_j-u_k}\le4\), the conditional variance is at most
\(16/d\).  Therefore
\[
 \Pp\left(
 \abs{(u_j-u_k)^\top Hv}\ge z\mid U
 \right)
 \le2e^{-dz^2/32}.
\]
A union bound over at most \(N^2\) ordered pairs gives
\[
 \Pp\left(w_U(Hv)\ge z\mid U\right)
 \le2N^2e^{-dz^2/32}.
\]
With
\[
 z^2=\frac{32}{d}(L_N+r),
 \qquad
 L_N=\log(2N^2),
\]
we obtain
\[
 \Pp\left(
 w_U(Hv)^2\ge\frac{32}{d}(L_N+r)
 \ \middle|\ U
 \right)
 \le e^{-r}.
\]

For \(i\in I\), define
\[
 Y_i
 :=
 \left(
 \frac d{32}w_U(Hv_i)^2-L_N
 \right)_+.
\]
Then, conditionally on \(U\), the variables \(Y_i\) are independent and
\[
 \Pp(Y_i\ge r\mid U)\le e^{-r}
 \qquad (r\ge0).
\]
For \(0<\lambda<1\), the tail-integration identity gives
\begin{align*}
 \E[e^{\lambda Y_i}\mid U]
 &=
 1+\int_0^\infty \lambda e^{\lambda r}
 \Pp(Y_i\ge r\mid U)\,dr\\
 &\le
 1+\lambda\int_0^\infty e^{-(1-\lambda)r}\,dr
 =\frac1{1-\lambda}.
\end{align*}
Thus, for \(b>1\), Chernoff's inequality with \(\lambda=1-1/b\) yields
\[
 \Pp\left(
 \frac1s\sum_{i\in I}Y_i\ge b
 \ \middle|\ U
 \right)
 \le
 e^{-s(b-1-\log b)}.
\]
For \(b\ge8\),
\[
 b-1-\log b\ge\frac b2.
\]
We claim that
\[
 s\bigl(b_s(\tau)-1-\log b_s(\tau)\bigr)\ge T(\tau).
\]
Indeed, if \(b_s(\tau)=2T(\tau)/s\), the claim follows from the preceding
half-linear bound.  If \(b_s(\tau)=8\), then \(T(\tau)\le4s\), while
\(8-1-\log8>4\).  This proves the claim.

Take a \(d^{-1/2}\)-net \(\mathcal H\) of the operator unit ball.  By
Lemma~\ref{lem:nets},
\[
 \abs{\mathcal H}\le(3d)^{d^2}.
\]
A union bound over \(H_0\in\mathcal H\) shows that, with conditional failure
probability at most
\[
 (3d)^{d^2}e^{-T(\tau)}=e^{-\tau},
\]
one has simultaneously
\[
 \frac1s\sum_{i\in I}w_U(H_0v_i)^2
 \le
 \frac{32}{d}\left(L_N+b_s(\tau)\right)
 \qquad (H_0\in\mathcal H).
\]
This defines \(\mathcal E_I(\tau)\).

Now let \(\norm{H}_{\op}\le1\), and choose \(H_0\in\mathcal H\) satisfying
\(\norm{H-H_0}_{\op}\le d^{-1/2}\).  On the input-norm event in the
statement,
\begin{align*}
 w_U(Hv_i)
 &\le
 w_U(H_0v_i)
 +\max_{j,k}\abs{(u_j-u_k)^\top(H-H_0)v_i}\\
 &\le
 w_U(H_0v_i)
 +4d^{-1/2}\norm{v_i}\\
 &\le
 w_U(H_0v_i)+\frac8{\sqrt d}.
\end{align*}
Therefore
\[
 w_U(Hv_i)^2
 \le
 2w_U(H_0v_i)^2+\frac{128}{d}.
\]
Average over \(i\in I\) and use the net estimate to obtain
\[
 \frac1s\sum_{i\in I}w_U(Hv_i)^2
 \le
 \frac{64}{d}\left(L_N+b_s(\tau)+2\right).
\]
Taking the supremum over \(H\) proves the result.
\end{proof}

\section{Operator-norm upper bound}
\label{sec:operator-upper}

Define the global norm event
\[
 \mathcal E_{\mathrm{norm}}
 :=
 \left\{
 \max_{j\in[N]}\norm{u_j}\le2,
 \quad
 \max_{i\in[N]}\norm{v_i}\le2
 \right\}.
\]
By Lemma~\ref{lem:gaussian-vector} and a union bound,
\begin{equation}
 \Pp(\mathcal E_{\mathrm{norm}})
 \ge1-2Ne^{-d/2}.
 \label{eq:norm-event-probability}
\end{equation}

Assume \(N>d\).  For \(r=0,\ldots,J-1\), set
\[
 a_r:=2^r d,
 \qquad
 I_r:=\{a_r+1,\ldots,\min(2a_r,N)\},
 \qquad
 s_r:=\abs{I_r}.
\]
These nonempty, disjoint blocks partition \(\{d+1,\ldots,N\}\).  For each
block, apply Lemma~\ref{lem:block-width} with \(\tau=d\).  Since
\(T(d)=T_d\), conditionally on any output realization satisfying the output
part of \(\mathcal E_{\mathrm{norm}}\), the simultaneous block event
\[
 \mathcal E_{\mathrm{width}}
 :=
 \bigcap_{r=0}^{J-1}\mathcal E_{I_r}(d)
\]
satisfies
\begin{equation}
 \Pp\left(\mathcal E_{\mathrm{width}}^c\mid U\right)
 \le J e^{-d}.
 \label{eq:width-event-probability}
\end{equation}
When \(N\le d\), define \(\mathcal E_{\mathrm{width}}\) to be the whole
sample space.

\begin{proposition}[Operator upper estimate on explicit events]
\label{prop:operator-upper}
On \(\mathcal E_{\mathrm{norm}}\),
\[
 \beta_{\op}\le16.
\]
If \(N>d\), then on
\(\mathcal E_{\mathrm{norm}}\cap\mathcal E_{\mathrm{width}}\),
\[
 \beta_{\op}
 \le
 16P_d
 +\frac{2^{\alpha+4}(L_N+10)}{d}
 +\frac{32T_d}{d^2}\,\Xi_{\alpha,J}P_d.
\]
\end{proposition}

\begin{proof}
On \(\mathcal E_{\mathrm{norm}}\),
\[
 R_u\le4,
 \qquad
 \Tr(M_v)=\sum_{i=1}^N p_i\norm{v_i}^2\le4.
\]
The second estimate in Lemma~\ref{lem:operator-reduction} gives
\[
 \beta_{\op}\le16.
\]

Assume now that \(N>d\).  For \(i\le d\) and
\(\norm{H}_{\op}\le1\),
\[
 w_U(Hv_i)
 \le
 \max_{j,k}\norm{u_j-u_k}\norm{Hv_i}
 \le4\cdot2=8.
\]
Thus the head contribution in Lemma~\ref{lem:operator-reduction} is at most
\[
 \frac14\sum_{i=1}^d p_iw_U(Hv_i)^2
 \le16P_d.
\]

Fix a tail block \(I_r\).  Since \(p_i\) is nonincreasing,
\[
 \sum_{i\in I_r}p_iw_U(Hv_i)^2
 \le
 p_{a_r+1}\sum_{i\in I_r}w_U(Hv_i)^2.
\]
The block-width estimate with \(\tau=d\) yields
\begin{align*}
 \frac14\sum_{i\in I_r}p_iw_U(Hv_i)^2
 &\le
 \frac{16}{d}p_{a_r+1}s_r
 \left[
 L_N+\max\left\{8,\frac{2T_d}{s_r}\right\}+2
 \right]\\
 &\le
 \frac{16(L_N+10)}d\,s_rp_{a_r+1}
 +\frac{32T_d}{d}\,p_{a_r+1}.
\end{align*}
For every \(i\in I_r\), we have \(i\le2a_r\), and therefore
\[
 p_i
 =\frac{i^{-\alpha}}{H_{N,\alpha}}
 \ge
 2^{-\alpha}\frac{a_r^{-\alpha}}{H_{N,\alpha}}
 \ge
 2^{-\alpha}p_{a_r+1}.
\]
Hence
\[
 s_rp_{a_r+1}
 \le
 2^\alpha\sum_{i\in I_r}p_i,
\]
and summing over blocks gives
\[
 \sum_{r=0}^{J-1}s_rp_{a_r+1}\le2^\alpha.
\]
Also,
\[
 p_{a_r+1}
 \le
 \frac{a_r^{-\alpha}}{H_{N,\alpha}}
 =p_d2^{-\alpha r}.
\]
Since \(P_d\ge dp_d\),
\[
 \sum_{r=0}^{J-1}p_{a_r+1}
 \le
 p_d\Xi_{\alpha,J}
 \le
 \frac{P_d}{d}\Xi_{\alpha,J}.
\]
Summing the preceding block estimate over \(r\) gives
\[
 \frac14
 \sup_{\norm{H}_{\op}\le1}
 \sum_{i=d+1}^N p_iw_U(Hv_i)^2
 \le
 \frac{2^{\alpha+4}(L_N+10)}d
 +\frac{32T_d}{d^2}\Xi_{\alpha,J}P_d.
\]
Adding the head contribution and invoking
Lemma~\ref{lem:operator-reduction} proves the claim.
\end{proof}

\section{Proof of the main theorem}
\label{sec:main-proof}

\begin{proof}[Proof of Theorem~\ref{thm:main}]
Because \(m\ge512\), we have \(K\ge1\), \(K\le d/256\), and \(2K\le N\).
Thus all selector and rectangular-matrix constructions above are valid.

Consider the event
\[
 \mathcal G
 :=
 \mathcal E_{\mathrm{norm}}
 \cap\mathcal E_M
 \cap\mathcal E_{\mathrm{tr}}
 \cap\mathcal E_{\mathrm{sel}}
 \cap\mathcal E_V
 \cap\mathcal E_G
 \cap\mathcal E_{\mathrm{width}}.
\]
The individual failure bounds are:
\begin{align*}
 \Pp(\mathcal E_{\mathrm{norm}}^c)
 &\le2Ne^{-d/2},\\
 \Pp(\mathcal E_M^c)
 &\le e^{-d},\\
 \Pp(\mathcal E_{\mathrm{tr}}^c)
 &\le e^{-d/16},\\
 \Pp(\mathcal E_{\mathrm{sel}}^c)
 &\le
 4K e^{-d/32}
 +2K e^{-d/48}
 +K(N-2)e^{-d/64},\\
 \Pp(\mathcal E_V^c)+\Pp(\mathcal E_G^c)
 &\le4\cdot17^K e^{-d/32}.
\end{align*}
For the width event, let
\[
 \mathcal E_{\mathrm{norm}}^U
 :=
 \left\{\max_j\norm{u_j}\le2\right\}.
\]
Equation~\eqref{eq:width-event-probability} and the tower property imply
\[
 \Pp\left(
 \mathcal E_{\mathrm{norm}}^U
 \cap\mathcal E_{\mathrm{width}}^c
 \right)
 \le J e^{-d}.
\]
Therefore a union bound gives
\[
 \Pp(\mathcal G^c)
 \le
 \varepsilon_{\mathrm{tot}}(d,N),
\]
with \(\varepsilon_{\mathrm{tot}}\) defined in
\eqref{eq:total-failure}.

On \(\mathcal G\), Proposition~\ref{prop:frobenius-upper} gives the
Frobenius upper bound, and Proposition~\ref{prop:frobenius-lower} gives the
Frobenius lower bound.  Proposition~\ref{prop:operator-lower} gives the
operator lower bound.  Proposition~\ref{prop:operator-upper} gives both the
universal operator upper bound \(16\) and, when \(N>d\), the sharper dyadic
upper bound.  Hence all conclusions of Theorem~\ref{thm:main} hold
simultaneously on \(\mathcal G\).
\end{proof}

\section{Zipf consequences in the polynomial-size regime}
\label{sec:consequences}

This section records consequences of Theorem~\ref{thm:main}; it is not used in
its proof.

\begin{corollary}[Polynomial-size summary]
\label{cor:polynomial-size}
Fix \(\alpha\ge0\).  Suppose \(N=d^{\Theta(1)}\) and
\(m=\min\{d,N\}\ge512\).  Then
\[
 \varepsilon_{\mathrm{tot}}(d,N)=e^{-\Omega(d)}.
\]
On the event in Theorem~\ref{thm:main},
\[
 \beta_F
 =\Theta\left(\frac1d+p_1\right),
\]
with universal implicit constants, and
\[
 \beta_{\op}
 =\widetilde\Theta_\alpha\left(\frac1d+P_m\right),
\]
where the tilde hides factors polylogarithmic in \(d\).
\end{corollary}

\begin{proof}
Since \(N\) is polynomial in \(d\), all polynomial prefactors in
\eqref{eq:total-failure} are dominated by the displayed exponential tails.
For the rectangular terms, use \(K\le d/256\):
\[
 17^K e^{-d/32}
 \le
 \exp\left[-\left(\frac1{32}-\frac{\log17}{256}\right)d\right]
 =e^{-\Omega(d)}.
\]
Thus \(\varepsilon_{\mathrm{tot}}(d,N)=e^{-\Omega(d)}\).

The Frobenius assertion is exactly Theorem~\ref{thm:main}.  For the operator
upper bound, if \(N\le d\), then \(P_m=P_N=1\), and the lower and upper bounds
are both of constant order.  If \(N>d\), then \(m=d\),
\[
 L_N=O(\log d),
 \qquad
 \frac{T_d}{d^2}=O(\log d),
 \qquad
 J=O(\log d),
\]
and
\[
 \Xi_{\alpha,J}
 =
 \begin{cases}
 O(\log d),&\alpha=0,\\
 O_\alpha(1),&\alpha>0.
 \end{cases}
\]
Therefore
\[
 \beta_{\op}
 \lesssim_\alpha
 \begin{cases}
 P_d\log^2d+\dfrac{\log d}{d},&\alpha=0,\\[2mm]
 P_d\log d+\dfrac{\log d}{d},&\alpha>0.
 \end{cases}
\]
Combining with the lower bound proves the stated polylogarithmic equivalence.
\end{proof}

\begin{corollary}[Evaluation of the Zipf head masses]
\label{cor:zipf-evaluation}
For fixed \(\alpha\ge0\),
\[
 H_{n,\alpha}
 =
 \begin{cases}
 \Theta_\alpha(n^{1-\alpha}),&0\le\alpha<1,\\
 \Theta(\log n),&\alpha=1,\\
 \Theta_\alpha(1),&\alpha>1.
 \end{cases}
\]
Consequently,
\[
 \beta_F
 =
 \begin{cases}
 \Theta_\alpha\!\left(d^{-1}+N^{\alpha-1}\right),&0\le\alpha<1,\\[1mm]
 \Theta\!\left(d^{-1}+(\log N)^{-1}\right),&\alpha=1,\\[1mm]
 \Theta_\alpha(1),&\alpha>1,
 \end{cases}
\]
and, in the polynomial-size regime,
\[
 \beta_{\op}
 =
 \begin{cases}
 \widetilde\Theta_\alpha\!\left(
 d^{-1}+\left(\dfrac{m}{N}\right)^{1-\alpha}
 \right),&0\le\alpha<1,\\[4mm]
 \widetilde\Theta\!\left(
 d^{-1}+\dfrac{\log m}{\log N}
 \right),&\alpha=1,\\[4mm]
 \Theta_\alpha(1),&\alpha>1,
 \end{cases}
 \qquad m=\min\{d,N\}.
\]
For \(0\le\alpha<1\) and \(N>d\), the exact identity is
\[
 P_d=H_{d,\alpha}p_1,
\]
and therefore
\[
 P_d
 \asymp_\alpha
 \left(\frac dN\right)^{1-\alpha}
 \asymp_\alpha
 d^{1-\alpha}p_1.
\]
The last two relations are asymptotic comparisons, not exact equalities in
general.
\end{corollary}

\begin{proof}
The harmonic-number estimates follow from elementary integral comparison.
Substitute
\[
 p_1=H_{N,\alpha}^{-1},
 \qquad
 P_m=\frac{H_{m,\alpha}}{H_{N,\alpha}}
\]
into Theorem~\ref{thm:main} and Corollary~\ref{cor:polynomial-size}.  The
identity \(P_d=H_{d,\alpha}p_1\) follows directly from the definitions.
\end{proof}

\end{document}
